\chapter{Bitácora de ceremonias}

Esta bitácora se reconstruye a partir de lo \textbf{verificable}: el historial de Git (175 commits, del 9 de septiembre al 20 de noviembre de 2026),
el release plan, el backlog y los ADR. Los sprints son \emph{lógicos}, definidos por la capacidad de producto que dejan funcionando; las
fechas son las de los commits que construyeron cada incremento. Las actas de las reuniones (planning, review y retrospectiva) no se versionan en el
repositorio: lo que cada ceremonia produjo y dejó escrito sí está, y es lo que se enlaza a continuación.

\section{Sprints, hitos y evidencia}
\begin{table}[H]
\centering\small
\caption{Sprints del proyecto, del Sprint~0 al 10, con su incremento y la evidencia que lo respalda.}\label{tab:sprints}
\begin{tabularx}{\textwidth}{llYY}
\toprule
\textbf{Sprint} & \textbf{Fechas} & \textbf{Incremento demostrable} & \textbf{Evidencia} \\
\midrule
0 & 9--10 sep & Estructura del repositorio (reactor Maven) y primeras pruebas de captura de imagen & Commits «Estructura inicial» e «inicio e intentos de captura» \\
\multicolumn{4}{l}{\textit{Release 1}} \\
1 & 10--19 sep & Detección de personas y maletas en el video (HU-01, 02) & «Detección exitosa de maleta y persona»; \cod{ObjectDetectorNmsTest} \\
2 & 19--20 sep & Registro de objetos en reposo (3~s quietos) con foto (HU-03, 04) & \cod{ObjectTrackerTest}; foto del frame en cada evento \\
3 & 20--21 sep & Retiro, dueño probable y tipo de retiro (HU-05, 06, 07, 09) & \cod{RemovalVerificationTest}, \cod{PersonDescriberTest} \\
4 & 21 sep & Alertas con evidencia y reglas de severidad (HU-13, 14, 18); descripción por IA opcional & \cod{AlertRulesTest}, \cod{AlertEngineTest} \\
5 & 21 sep--1 oct & Panel en tiempo real con WebSocket y alarma sonora; Docker (HU-15, 16, 17, 36) & \cod{AlertGatewayTest}; prueba con cámara real; «Ya funciona el Docker» (1 oct) \\
\multicolumn{4}{l}{\textit{Release 2}} \\
6 & 10--11 oct & Ventana compartida desde el navegador; dos cámaras con zonas; mediciones de rendimiento (HU-21 a 27, 37) & \cod{BrowserFrameSourceTest}, \cod{MultiCameraTest}; ADR~0003 \\
7 & 11 oct & Reporte en tabla y pista «dueño que volvió» (HU-12, 19, 38) & \cod{OwnerReturnTest}, \cod{ReportJsonTest} \\
8 & 11--17 oct & Avisos de texto (navegador y Telegram) y \textbf{inicio de sesión con Google} (12 oct); pipeline de CI (13 oct) & Servicio de notificaciones; mensaje de prueba recibido en Telegram; ADR~0004 \\
9 & 28 oct--12 nov & \textbf{Observabilidad} (Micrometer, Prometheus y Grafana, 28--29 oct) y \textbf{persistencia} (PostgreSQL, MongoDB y Redis; servicios de acceso y colaboración; 6--7 nov); tablero de Grafana del equipo & ADR~0005; evidencias numeradas 01 a 23; tablero exportado \\
10 & 13--22 nov & \textbf{Despliegue en AWS} con Caddy y dominio (13 nov), despliegue por SSH (16 nov); \textbf{seguridad} ES-1 a ES-3 (19 nov); \textbf{disponibilidad y escalado} (19 nov); \textbf{rendimiento} con k6 y \textbf{mantenibilidad} con Sonar (19--20 nov); este documento y el artículo & ADR~0006; informes de \cod{docs/}; 479 pruebas Java y 144 de Node \\
\bottomrule
\end{tabularx}
\end{table}

\section{Planning, Review y Retrospectivas}
\begin{table}[H]
\centering\small
\caption{Lo que cada ceremonia dejó escrito en el repositorio.}\label{tab:ceremonias}
\begin{tabularx}{\textwidth}{lYY}
\toprule
\textbf{Ceremonia} & \textbf{Artefacto que produjo} & \textbf{Dónde} \\
\midrule
Planning (qué entra al sprint) & Backlog con 38 HU en Gherkin, DoR y DoD, MoSCoW y story map; release plan por resultado de producto & \cod{docs/backlog.md}, \cod{docs/release-plan.md} \\
Decisiones de arquitectura & Seis ADR con alternativas, consecuencias y evidencia & \cod{docs/adr/} \\
Review (qué se demostró) & Verificación con cámara real y navegador; pruebas automáticas por incremento; guía y rúbrica de la sustentación & \cod{docs/sustentacion.md}; evidencias de \cod{docs/persistencia/} \\
Retrospectiva (qué se corrigió) & Cierre del feedback de la revisión interna, con decisión, evidencia y estado por observación & \cod{docs/sustentacion.md}, sección 5 \\
\bottomrule
\end{tabularx}
\end{table}

\section{Lecciones aprendidas}
Cada una salió de un hecho del proyecto, no de una recomendación general:

\begin{enumerate}
  \item \textbf{Medir antes de suponer.} «Menos de 5 segundos» nunca se había medido: \cod{AlertLatencyBenchmark} dio media 0,5~s y P95 0,9~s. Y las
  pruebas de k6 mostraron que la CPU no era el límite: era una conexión única con consultas repetidas (de 34--48 a 570--678 req/s al corregirla).
  \item \textbf{El análisis estático y las pruebas encuentran lo que el ojo no ve.} Medir la cobertura del panel destapó dos defectos reales;
  Sonar encontró registros con arreglos que comparaban identidad. Y un arreglo sin prueba hizo fallar el control de calidad (el código nuevo cuenta).
  \item \textbf{Los servicios externos mueren.} El modelo de visión de NVIDIA dejó de existir (HTTP 410) y la descripción del dueño salió vacía hasta
  cambiarlo; por eso todo canal externo es opcional, en segundo plano y no frena la alerta.
  \item \textbf{La caché puede mostrar, no autorizar.} El laboratorio de consistencia lo probó: con el permiso cacheado, una edición pasaba aunque la
  autoridad estuviera caída; ahora se rechaza con 503.
  \item \textbf{Declarar los límites sube la credibilidad.} El despliegue es de un solo servidor; la prueba con dos cámaras reales y el correo siguen
  pendientes; los resultados de rendimiento son de un equipo con el generador incluido. Todo está dicho en el texto.
  \item \textbf{El entorno también falla.} Windows suspendió el equipo en pleno estrés y dos corridas quedaron inválidas; el script ahora pide no
  suspender y avisa si la prueba tardó de más.
  \item \textbf{Los diagramas se verifican contra el código.} Las vistas del diseño inicial mostraban componentes que nunca existieron (un bus de
  eventos, una base de alertas separada); se rehicieron desde el código.
\end{enumerate}

\section{Contribuciones individuales}\label{sec:contribuciones}
La Tabla~\ref{tab:contribuciones} se obtiene de \cod{git log} y es verificable; el número de commits no mide el esfuerzo, y el trabajo de
revisión, pruebas con cámara real y decisiones no siempre deja commits.

\begin{table}[H]
\centering\small
\caption{Contribuciones por persona (commits y áreas, de git log).}\label{tab:contribuciones}
\begin{tabularx}{\textwidth}{lrY}
\toprule
\textbf{Integrante} & \textbf{Commits} & \textbf{Áreas principales} \\
\midrule
Hever Barrera Batero & 130 & Núcleo de detección, seguimiento y alertas; gateway y panel; persistencia y servicios de acceso y colaboración; seguridad, disponibilidad, rendimiento y mantenibilidad del Release~2; ADR y este documento \\
María Juliana Rodríguez Caicedo & 29 & Inicio de sesión con Google; pipeline de CI/CD y despliegue por SSH; calidad con Sonar y JaCoCo; despliegue en AWS (Caddy, HTTPS, dominio, reinicio automático); protección de Grafana; pruebas de calidad \\
Kevyn Daniel Forero Gonzalez & 16 & Estructura inicial del repositorio; integración de back y front con Docker; instrumentación con Micrometer y métricas de alertas; Prometheus, Grafana y los tableros \\
\bottomrule
\end{tabularx}
\end{table}

\section{Uso de inteligencia artificial}
\begin{table}[H]
\centering\small
\caption{Declaración de uso de IA (descripción, no conversaciones completas ni información confidencial).}\label{tab:ia}
\begin{tabularx}{\textwidth}{lYY}
\toprule
\textbf{Herramienta} & \textbf{Para qué se usó} & \textbf{Cómo se validó} \\
\midrule
YOLOv8 (local) & Detectar personas y objetos y estimar la pose del dueño (parte del producto) & Pruebas automáticas, calibración con fotos reales, pruebas con cámara \\
Modelo de visión de NVIDIA (servicio externo, opcional) & Describir al dueño y a quien retira el objeto (parte del producto) & Pruebas contra un servidor falso; el recorte con la cara sale a un tercero, por eso es opcional y declarado \\
Claude Code (asistente de desarrollo, Anthropic) & Apoyo en análisis, código, pruebas, scripts de k6, workflows de CI, informes, diagramas y documentación (ADR, backlog, este documento) & El equipo revisó el código, ejecutó las pruebas (479 de Java y 144 de Node), midió con k6 y Sonar, probó con cámara real y decidió; la responsabilidad del resultado es del equipo \\
\bottomrule
\end{tabularx}
\end{table}

Ningún secreto, credencial ni dato sensible se incluyó en las instrucciones a la IA ni en el repositorio; las claves van solo en variables de
entorno. Los integrantes deben poder explicar y defender el código y las decisiones, incluidos los asistidos por IA.
